% Copyright (c) 2026 Geovana Neves. Licensed under CC BY 4.0 (https://creativecommons.org/licenses/by/4.0/). See LICENSE-AND-AUTHORSHIP.md.
%
% 01-workspaces/text/04-running.tex

\section{Running}

\begin{frame}{Four stages}
\slidefig[0.50]{guide-stage-flow}
{\ftstablesize
\begin{tabular}{@{}lll@{}}
\toprule
\textbf{Stage} & \textbf{Windows} & \textbf{Cluster} \\
\midrule
\pyfs{plan}    & required, no seat & required, no queue \\
\pyfs{run}     & solves, then returns & submits, returns at once, \code{SUBMITTED} \\
\pyfs{collect} & not needed & waits for settled outputs, judges them \\
\pyfs{post}    & rebuilds the products & the same \\
\bottomrule
\end{tabular}}
\end{frame}

\begin{frame}{One matrix, any build}
\relax
\small
\begin{tabular}{@{}lllll@{}}
\toprule
 & 25.000 & 25.100, 26.000 & 26.100 & 26.101 and later \\
\midrule
\code{steady} & refused & runs & runs & runs \\
\code{unsteady} & refused & single march & single march & single march, or actions \\
\code{unsteady\_rotor} & refused & Euclidean rotor & unmarked, rev/min & rotary, or actions \\
\bottomrule
\end{tabular}

\vspace{2mm}
\begin{itemize}
  \item Only \code{FS\_BUILD} changes. Actions exist from 26.122.
  \item Without actions: \code{WALLTIME}, the snapshot thresholds and
  \code{FINISH\_PENDING}/\code{ADDITIONAL\_REVS} are \textbf{refused by name} at plan.
  \item \code{march\_strategy} in \code{plan.json}, the record and the superfile.
  \item A preset or pproc file can still block an older build: \textbf{plan first}.
\end{itemize}
\end{frame}

\begin{frame}{Tip and helical Mach, on every rotor point}
\[ \Omega = \frac{2\pi\,\mathrm{RPM}}{60}, \qquad
   M_\mathrm{tip} = \frac{\Omega R}{a}, \qquad
   M_\mathrm{hel} = \frac{\sqrt{V^2 + (\Omega R)^2}}{a} \]
\begin{itemize}\small
  \item For each rotor an \code{unsteady\_rotor} or \code{qsteady\_rotor} row
        turns, a \code{steady} row that states \code{RPM}, and each actuator
        disc a row names. \(R\)
        is the rotor's half diameter or the disc's tip radius; \(V\) and
        \(a\) are the point's own resolved free stream and speed of sound.
  \item \pyfs{plan} prints both per rotor and per point, and \textbf{warns},
        naming the point, when \(M_\mathrm{hel} \ge 1\): the tip is sonic or
        supersonic, where the linearised compressibility of a panel method
        stops holding [4, 12, 13]. A warning, never a refusal.
  \item The rotor table ends with \code{MTIP\_<alias>} and
        \code{MHEL\_<alias>}. Both are geometric: \(V\) is the free stream
        alone, without the rotor's own induced velocity.
\end{itemize}
\end{frame}

\begin{frame}{\pyfs{run}, on each machine}
\begin{columns}[T]
\begin{column}{0.48\textwidth}
\begin{block}{Windows}
\begin{itemize}\small
  \item A steady angle sweep is \textbf{one job over all its points}; each
        point starts cold. \code{COLD\_START: false} starts each point from
        the previous one.
  \item \code{WALLTIME} arms a clock; a run it stopped is
        \code{WALLTIME\_REACHED}, \textbf{not a failure}.
\end{itemize}
\end{block}
\end{column}
\begin{column}{0.48\textwidth}
\begin{block}{Cluster}
\begin{itemize}\small
  \item An unsteady point is \textbf{its own job}, in its own datapoint folder
        with its descriptor, action program and clock.
  \item A steady angle sweep is one job; a flow-variable sweep is one per point.
  \item A point already queued is refused a second submission. Collect it
        first.
\end{itemize}
\end{block}
\end{column}
\end{columns}
\end{frame}

\begin{frame}[fragile]{\pyfs{collect} watches the workspace, not the scheduler}
\begin{lstlisting}[style=ftsbase]
pyfs-matrix collect --workspace .            # one sweep, what a cron wants
pyfs-matrix collect --workspace . --watch    # loop until nothing is outstanding
*/15 * * * *  cd <study> && $FTS/py-venv/(*@\envname@*)/bin/pyfs-matrix collect --workspace .
\end{lstlisting}
\begin{columns}[T]
\begin{column}{0.48\textwidth}
\begin{block}{Settled is two signals}
\small Size and time stable across two looks, \textbf{and} the log present,
because the log is declared last.
\end{block}
\end{column}
\begin{column}{0.48\textwidth}
{\ftstablesize
\begin{tabular}{@{}cl@{}}
\toprule
\textbf{Exit} & \textbf{Meaning} \\
\midrule
0 & everything settled and collected \\
1 & a point it could not complete \\
3 & still outstanding \\
\bottomrule
\end{tabular}}
\end{column}
\end{columns}
\end{frame}

\begin{frame}{\code{RESTART}: continuing a march that stopped}
% \relax first: a brace group right after the title is read as a subtitle.
\relax
{\ftstablesize
\begin{tabular}{@{}ll@{}}
\toprule
\textbf{Write in the free cell} & \textbf{When} \\
\midrule
\code{RESTART: \{FINISH\_PENDING\}}     & a run the \emph{clock} stopped \\
\code{RESTART: \{ADDITIONAL\_ITERS=n\}} & any continuable run, $n$ more steps \\
\code{RESTART: \{ADDITIONAL\_REVS=n\}}  & any continuable run, $n$ more revolutions \\
\bottomrule
\end{tabular}}
\vspace{3mm}
\begin{itemize}
  \item You write \code{RESTART} and nothing else; the manifest says which run
        and how many steps are left. The count is a \textbf{remainder}.
  \item It continues the point's \textbf{most recent} run. A finished point is
        refused; a failed continuation is not retried.
  \item What it replaces is archived first, under a day and hour stamp.
\end{itemize}
\begin{takeaway}On the cluster this is the ordinary case: the queue enforces
the wall clock, so \code{\{FINISH\_PENDING\}} is the form you will use most.\end{takeaway}
\end{frame}

\begin{frame}[fragile]{Points already run, and a Linux run kept local}
\begin{lstlisting}[style=ftsbase]
pyfs-matrix run matriz.fs --workspace . --resume
pyfs-matrix run matriz.fs --workspace . --force-rerun '<point name or run id>'
pyfs-matrix run matriz.fs --workspace . --local
\end{lstlisting}
\begin{itemize}\small
  \item A recorded point is refused before anything runs; \pyfs{plan} marks it
        \code{ALREADY\_RECORDED}.
  \item \code{--resume}: only the new points of a matrix that grew.
  \item \code{--force-rerun}: redo a named point; its old record and outputs
        move to \code{archive/}, nothing is deleted. Plan again first.
        Not together with \code{--resume}.
  \item \code{--local}: on Linux with a profile, run here instead of
        submitting; the record says \code{forced\_local}. Elsewhere it changes
        nothing.
\end{itemize}
\begin{alertblock}{Redo on purpose}
\small A recorded point is never solved again by \code{--resume}. A point
whose inputs were wrong is redone by name with \code{--force-rerun}.
\end{alertblock}
\end{frame}

\begin{frame}[fragile]{The additional post}
\begin{lstlisting}[style=ftsbase]
ADDITIONAL_PPROC: p002                        # free cell; the run is unchanged
pyfs-matrix post matriz.fs --workspace . --additional-pproc
\end{lstlisting}
\begin{columns}[T]
\begin{column}{0.50\textwidth}
{\ftstablesize
\begin{tabular}{@{}ll@{}}
\toprule
\textbf{Reopened \code{.fsm}, no solve} & \textbf{Measured [3]} \\
\midrule
total loads, surface, sections & identical \\
sectional loads & identical once computed \\
unsteady plots history & identical \\
probe points off the body & \alert{not identical} \\
\bottomrule
\end{tabular}}
\end{column}
\begin{column}{0.46\textwidth}
\begin{itemize}\small
  \item Opens a copy, cuts the sections, exports, closes. Never solves or saves.
  \item \code{p002} may ask sections, their loads and the surface only.
  \item 26.124 only. Unsteady: the last instant.
  \item On the cluster: submitted, recorded \code{SUBMITTED};
        \pyfs{collect} records \code{EXTRACTED}.
\end{itemize}
\end{column}
\end{columns}
\end{frame}

\section{What a run leaves}

\begin{frame}{A simulation folder}
\begin{columns}[c]
\begin{column}{0.58\textwidth}
\slidefig[0.88]{guide-sims-anatomy}
\end{column}
\begin{column}{0.40\textwidth}
\begin{itemize}\small
  \item \code{inputs/}: a link to the mesh's library folder.
  \item \code{scripts/}: what the solver was handed. The fastest answer to
        \emph{what did this run really do}.
  \item \code{profiles/}: where the probes were placed.
  \item \code{datapoints/DP-<tag>/}: one folder per point.
\end{itemize}
\begin{takeaway}Nothing here is edited. Everything here is reproducible from
the inputs.\end{takeaway}
\end{column}
\end{columns}
\end{frame}

\begin{frame}{Inside a datapoint: what can be deleted}
% \relax first: a brace group right after the title is read as a subtitle.
\relax
{\ftstablesize
\begin{tabular}{@{}lll@{}}
\toprule
\textbf{File} & \textbf{Holds} & \textbf{Read afterwards by} \\
\midrule
\code{<stem>.fsm}         & saved simulation & \code{RESTART}. Delete it and the point cannot continue \\
\code{<stem>.txt}         & loads            & \code{polars/}, \code{campaign\_sweep.csv}, the judge \\
\code{<stem>\_sloads.txt} & sectional loads  & \code{sections/} \\
\code{<stem>\_probes.txt} & probe samples    & \code{probes/} \\
\code{<stem>\_plots.txt}  & per-step plots (unsteady) & polar, rotor averages, reductions \\
\code{<stem>\_cp.txt}, \code{<stem>.dat} & section pressures, Tecplot & nothing: a viewer \\
\code{<stem>\_log.txt}    & solver log       & \code{collect}; the first file to open \\
\bottomrule
\end{tabular}}
\vspace{2mm}
\small The stem is \code{P<POL>-<point name>}, for example
\code{P6112-M200AL+020BE+000}; the name follows the declared condition keys.
\end{frame}

\begin{frame}{Eight states, and four are not failures}
\begin{columns}[T]
\begin{column}{0.48\textwidth}
{\ftstablesize
\begin{tabular}{@{}ll@{}}
\toprule
\textbf{Not a failure} & \\
\midrule
\code{CONVERGED}           & reached its criterion \\
\code{COMPLETED\_MAX\_ITER} & reached its cap \\
\code{WALLTIME\_REACHED}    & clock stopped it \\
\code{SUBMITTED}           & in the queue \\
\bottomrule
\end{tabular}}
\end{column}
\begin{column}{0.48\textwidth}
{\ftstablesize
\begin{tabular}{@{}ll@{}}
\toprule
\textbf{Failure} & \\
\midrule
\code{FAILED\_DIVERGED}           & diverged \\
\code{FAILED\_EXECUTION}          & solver did not run \\
\code{FAILED\_SCRIPT}             & script rejected \\
\code{FAILED\_INCOMPLETE\_OUTPUT} & outputs missing \\
\bottomrule
\end{tabular}}
\end{column}
\end{columns}
\vspace{3mm}
\begin{takeaway}\code{WALLTIME\_REACHED} and \code{COMPLETED\_MAX\_ITER} carry real
numbers, and are read wrongly more often than any real failure.\end{takeaway}
\end{frame}

\section{What post-processing leaves}

\begin{frame}{Who writes the products}
\slidefig[0.42]{guide-post-anatomy}
{\ftstablesize
\begin{tabular}{@{}ll@{}}
\toprule
\pyfs{run}     & at its end, unless it submitted; no flag turns it off \\
\pyfs{collect} & after a sweep that collected something, and only then; \code{--no-post} \\
\pyfs{post}    & when you ask; no seat, no solver \\
\bottomrule
\end{tabular}}
\vspace{1mm}

\small All three \textbf{archive} what they replace, into
\code{archive/<day and hour>/}. Only \code{--force-overwrite} keeps no copy, and it
asks first. It is not part of a routine. \code{series/} is archived too.
\end{frame}

\begin{frame}{The five product folders}
% \relax first: a brace group right after the title is read as a subtitle.
\relax
{\ftstablesize
\begin{tabular}{@{}p{0.16\linewidth}p{0.50\linewidth}p{0.26\linewidth}@{}}
\toprule
\textbf{Folder} & \textbf{Answers} & \textbf{Plot against} \\
\midrule
\code{polars/}     & what the coefficients do as the condition changes; one file per group,
                     body, stability and wind axes & \code{ALPHA} or \code{BETA} \\
\code{series/}     & how it got there: loads, sections and probes per step & \code{azimuth\_deg} \\
\code{sections/}   & sectional loads, in each declared frame & \code{Offset} \\
\code{probes/}     & the flow samples, and time-average, phase-locked, per-blade reductions & \code{STEP}, \code{AZIMUTH}, \code{BLADE} \\
\code{provenance/} & inputs, script, outputs with sha256, argv, build & makes a plot defensible \\
\bottomrule
\end{tabular}}
\vspace{3mm}
\begin{itemize}\small
  \item A rotor row reduces \textbf{per rotor}: \code{<point>\_per\_blade\_<ALIAS>.csv}.
  \item In \code{campaign\_sweep.csv}, check \code{fs\_build} first: the build
        the solver reported.
\end{itemize}
\end{frame}

\begin{frame}{\code{additional/}, and what the products hold}
\begin{itemize}\small
  \item \code{post/<matrix>/additional/<pid>/}: the additional post's polars,
        sections and, unsteady, plots reductions. Every \code{products.json}
        entry says \code{"additional": true}, \code{pproc},
        \code{extraction}, \code{derives\_from}; filter on it.
  \item Its sections table keeps the run's rows first, then the new ones;
        \code{leading\_sections} counts the run's.
  \item An induced drag nobody computed is \code{NA}, with the axis columns it
        reaches. A steady probe line is exported once.
  \item Unsteady \code{series/<point>\_loads\_series.csv}: moments about the
        moment point.
  \item \code{post.log.json} beside \code{post.log}. Post a workspace with
        the package version that ran it.
\end{itemize}
\end{frame}

\begin{frame}{Vector loads and window averages}
\begin{itemize}\small
\item Steady polar axes use the exported force and moment vectors.
\code{CDB = Cx}, \code{CLB = Cz}; sideslip points get rows.
\item \code{CLS}/\code{CLW} fall by 0.10 to 0.25 per cent on 27 of 28 recorded
lifting exports, including zero sideslip; one differs by 0.71 per cent.
The cause of the export's lift discrepancy is unknown.
\item Unsteady rotor tables are window averages of all six global
\code{MRP} components over exactly the rotor's families. Missing history is a
skip, never a last-step fallback. A run adds \code{ROTOR\_<ALIAS>} if needed.
\item \code{RPM\_<alias>} and \code{DIAMETER\_<alias>} join the rotor header;
\code{J\_<alias>} uses the free stream. $\mathrm{ETAW}=J\,\mathrm{CTW}/\mathrm{CP}$
uses the whole force in wind axes and equals \code{ETA} along the stream.
\item \code{ETAW} and the shaft angle follow the stream direction at
incidence or sideslip; a re-post needs no solver.
\end{itemize}
\end{frame}

\begin{frame}{Read columns by name}
\begin{itemize}\small
\item Condition: \code{ALPHA, BETA, MACH, RE, VINF, VREF, ALT, RHO, TEMP, MU,}
\code{J, SREF, CREF, BREF}. Plots keep the export's own headings.
\item Unsteady polar: \code{P<sim>\_<name>\_uns\_avg.csv}, beginning with
\code{FIRST\_STEP, LAST\_STEP, STEPS}; adds moment point, group-suffixed axes,
equations and super-file content.
\item Axes need all six global \code{MRP} plots. A new run supplies
\code{MRP\_TOTAL} if absent; a history without them has a \code{\#axes} skip.
\item Reductions carry \code{ROTOR} and \code{XMOM, YMOM, ZMOM}, not
\code{Time-step} and \code{Time (sec)}.
\item Sections: \code{STEP, FAMILY, PLANE, ROTOR, AZIMUTH}; still one instant.
Azimuth is blade one of that row's rotor, or \code{NA} without one.
\item Series use \code{STEP}, add the condition block; probes add \code{PROBE}.
\end{itemize}
\end{frame}

\begin{frame}{One row per blade, or per azimuth}
\begin{itemize}\small
\item \code{per\_blade}: one row per blade, each blade's own rotor clock and matrix window,
\code{BLADE, FAMILY, AZIMUTH\_START, AZIMUTH\_END}. A blade's family suffix is
removed from its own plot columns so all blades share headings.
\item Declared \code{[phase\_locked]}: one row per azimuth of the last
revolution, averaged across the requested revolutions. Blade columns align by
each blade's own azimuth; other columns use blade one.
\item The row's total revolutions must be \emph{at least}
\code{min\_revolutions}. A shorter run loses only that table. Without the
pproc declaration, the file remains a passage series.
\item Editing the matrix window moves the reductions together at post;
editing the pproc gate or depth also needs no new solver run. Records stay as written.
\item Read \code{products.json}: \code{complete}, \code{interrupted}, sources,
windows and skips. \code{\#axes}, \code{\#equations}, \code{\#names} count for
\code{post --strict} even if the rest of the file was written.
\end{itemize}
\end{frame}

\begin{frame}{Collection and manifest limits}
\begin{itemize}
\item An early export at an iteration-limit run, confirmed by the solver log,
and a velocity mismatch are failures and leave the products.
\item A continuation's \code{continues} field excludes its predecessor from
products and the sweep table.
\item Stamped series exports are found in both simulation and datapoint
folders; a missing kind is a skip, not a header-only product.
\end{itemize}
\begin{alertblock}{Manifest ownership}
A live writer renews its owned lock every 5 seconds. Takeover requires a dead
local owner or a heartbeat older than 300 seconds. Finish every older writer
before a newer one uses the same workspace.
\end{alertblock}
\end{frame}

\begin{frame}{Reading a result honestly}
\begin{itemize}
  \item For unsteady, \code{CONVERGED} checks the last time step's inner
        iterations. Whether the history has settled is the user's judgement.
        It does not mean the mesh was
        fine, the reference right, or the physics the one you meant.
  \item A curve that stops early at \code{WALLTIME\_REACHED} is real up to that
        step. Continue it, do not repeat it.
  \item One diverged point in a swept row is judged on its own.
\end{itemize}
\begin{ftslimit}
The workspace makes a study reproducible and traceable. That is not the same
as correct, and no amount of green status is evidence of the second.
\end{ftslimit}
\end{frame}

\begin{frame}{Re-posting a recorded campaign}
\begin{itemize}\small
  \item Run \code{pyfs-matrix post --workspace <root>} without a solver.
        Records remain unchanged; each product is written or skipped with its reason.
  \item Nothing blocks by default: frozen-solve averages are
        WARNED in \code{post.log} and written; \code{--check-frozen} withholds
        them. The unsteady polar's \code{ADVANCE\_RATIO} states the point's \code{J}.
  \item Every unsteady probe uses fluid plots. A cited profile recorded only
        as an instant needs a new run for its missing history.
  \item Each distribution gains \code{\_sloads\_<name>.csv} and
        \code{\_cp\_<name>.csv}, with every available exported \code{STEP}.
  \item New VTK/CSV exports and the pproc surface-averaging window need a new
        solver run. Guides 4 and 5 have the keys [4].
  \item \code{integrate = true} on a sections entry adds the
        strip force and moment columns on re-post, without a solver run.
\end{itemize}
\end{frame}
